\section*{Bedeutung von Spielen}
\par\noindent\rule{\textwidth}{0.4pt} \newline \newline
Spiel ist nicht gleich Spiel. Es gibt verschiedene Arten von Spielen, die jeweils verschiedene Ziele 
verfolgen. Als Gruppenleiter muss man sich im Vorfeld der Spielauswahl überlegen, was mit dem 
Spiel bezweckt werden soll. Beispielsweise wäre ein Kennenlernspiel nicht das richtige, um 
gemeinsam in der Gruppe mal „Dampf“ abzulassen. Folgende Arten von Spielen können 
unterschiedenen werden:
\begin{itemize}
    \item Kennenlernspiele: Die Anfangssituation einer Gruppe gestalten, sich gegenseitig 
(besser) kennen lernen und Vertrauen zu den anderen finden. 
\item Spiele zur Gruppeneinteilung: Kreativ, sowie abwechslungsreich Gruppen bilden und 
dabei „in Fahrt“ kommen.
\item Kreisspiele: Die perfekten Spiele für das Warming-Up der Gruppenstunde oder für 
die actionreiche Unterbrechung der Einheit. 
\item Bewegungsspiele: Nicht rumsitzen, sondern zusammen im Gruppenraum oder 
draußen aktiv werden.
\item Geländespiele: Auf dem Bolzplatz oder im Wald mit der ganzen Gruppe ein kleines Abenteuer erleben.
\item Spiele für die Abendrunde: Gemeinsam in der Gruppe den Tag ausklingen lassen – egal ob leise oder laut.
\end{itemize}
\section*{Was muss beim Spielen beachtet werden? }
\par\noindent\rule{\textwidth}{0.4pt} \newline \newline
So vielfältig die Auswahl an Spielen ist, so vielfältig sind auch die Voraussetzungen, die jeweils 
erfüllt sein müssen. Vor jedem Spiel sollten daher folgende Fragen beantwortet werden:
\begin{itemize}
    \item Was möchte ich mit dem Spiel bezwecken?
    \item Passt das Spiel zum Alter und zur derzeitigen Phase der Gruppe?
    \item Welche (Sicherheits-)Risiken gibt es zu beachten?
    \item Welches Material brauche ich für das Spiel und was muss ich im Vorfeld vorbereiten? Wenn im Rahmen einer Spielekette Spiele aneinander gereiht werden, ist auf die richtige Anordnung zu achten. So sollten am Anfang leichte und gut verständliche Spiele stehen, die alle Teilnehmer verstehen und ohne Probleme mitspielen können. 
\end{itemize}
Im Folgenden können die Spiele dann bezüglich Schwierigkeitsgrad, Komplexität, Teamarbeit und 
Körperkontakt stetig zunehmen, so dass das letzte Spiel der Gruppe als Team am meisten 
abverlangt. Zum einen wird durch eine solche Anordnung Spannung aufgebaut und Abwechslung 
erreicht. Zum anderen hilft dieser Aufbau aber auch den Gruppenmitgliedern, sich besser auf die 
Spielekette und die Gruppe einzulassen.
\section*{Spielerklärung}
\par\noindent\rule{\textwidth}{0.4pt} \newline \newline
Es gibt sicherlich in einigen Gruppen Spiele, die wie von selbst laufen, da sie im Laufe der Zeit zu 
„Klassikern“ geworden sind, die keiner Erklärung mehr bedürfen. Im Normalfall steht vor einem 
Spiel jedoch eine Spielerklärung durch den Gruppenleiter. 
Diese Erklärung kann natürlich ganz ohne irgendwelche Ausschmückungen erfolgen, 
beispielsweise so:
\begin{itemize}
    \item[]„Jetzt spielen wir das Spiel ‚Kissenrennen‘. Zählt mal bitte Teams ab: 1-2-1-2-1-2 usw. Alle, die in Mannschaft 1 sind, gehören zum blauen Kissen. Alle aus Team 2 gehören zum roten. Wenn ich gleich ‚los‘ sage, gebt ihr jeweils euer Kissen an euern linken Teamnachbarn weiter. Wenn ein Kissen das andere überholt, habt ihr gewonnen! Habt ihr noch Fragen? Ok, dann auf die Plätze, fertig, los!“
\end{itemize}
    

\section*{Spielepädagogik}
\par\noindent\rule{\textwidth}{0.4pt} \newline \newline
Spannender wird es allerdings, wenn zu den Regeln eine Geschichte erzählt wird, die das Spiel in 
einen (erfundenen) Zusammenhang setzt und damit die Mitspieler motiviert:
\begin{itemize}
    \item[]"Wir befinden uns auf der alt-ehrwürdigen Pferderennbahn des Jugendhofes in 
Vechta. Hier werdet ihr in Kürze Zeugen eines einmaligen Spektakels des 
Pferdesports: Es kommt zum lang ersehnten Schlagabtausch der Spitzenrennpferde 
Horst und Wendy! Applaus bitte (Teilnehmer jubeln)! Horst (blaues Kissen 
hochhalten), ein alter, erfahrener Hengst, macht heute vielleicht das letzte Rennen 
seiner Karriere! Kann er es mit einem Sieg abschließen? Ihm gegenüber steht Wendy 
(rotes Kissen hochhalten), eine zarte Stute im bestem Rennalter, die bisher noch kein 
Turnier gewinnen konnte, aber hier vielleicht vor ihrem ersten Erfolg steht. Bevor 
das Rennen beginnt, müssen wir die Fanlager bestimmen: Horst-WendyHorst Wendy-Horst-Wendy usw. (abzählen). Heute ist es Aufgabe der Fans, also euch, euer 
jeweiliges Pferd zum Sieg zu verhelfen! Dazu sollt ihr es nicht nur anfeuern, sondern 
auch antreiben…“
\end{itemize}

Bei der Wahl der Ausschmückungen ist immer darauf zu achten, dass diese für die jeweilige 
Gruppe passend sind. Eine Gruppe 15-jähriger Jungen freut sich beispielsweise in der Regel nicht 
darüber, in „das Land der Feen“ entführt zu werden. \\

Die Art und Weise, wie man einzelne Spiele motivierend erklärt, gilt (besonders) auch für 
Spieleketten. Anstatt einzelne Spiele isoliert aneinander zu hängen, könnten sie beispielweise 
Stationen einer Safari durch den Dschungel sein, die man gemeinsam unternimmt. Hierbei stellen 
sich am Anfang viele Fragen, die die Teilnehmer in den Bann der Geschichte ziehen:

\begin{itemize}
    \item Wo genau findet die Safari statt? Wie sieht es dort aus (evtl. Requisiten verwenden)? 
    \item In welcher Zeit findet die Safari statt? Was ist besonders für die Zeit? 
    \item Welche Rollen haben die Teilnehmer auf der Safari und welche der Gruppenleiter?
    \item Was ist der Grund oder das Ziel der Safari? 
    \item Kennen die Teilnehmer Grund oder Ziel oder kommen sie erst während der Safari zu einer Erkenntnis?
\end{itemize}

Der Kreativität des Gruppenleiters sind hier keine Grenzen gesetzt, solange die jeweilige 
Geschichte für die Zielgruppe passend ist. Sofern die Spielekette in einem zeitlich begrenzten 
Rahmen stattfindet (z.B. in einer Gruppenstunde), muss darauf geachtet werden, dass das Ziel der 
Geschichte erreicht wird. Daher lieber einen variablen Mittelteil einplanen, der bei Zeitknappheit 
gekürzt werden kann.
